% =====================================================================
% Section file for Overleaf: AI-READI Counterfactual Strategy
% Suggested placement: sections/08_counterfactual_strategy.tex
% Required packages in main.tex if not already present:
% \usepackage{amsmath,amssymb,booktabs}
% Optional: \usepackage{enumitem}
% Bibliography keys are provided in aireadi_counterfactual_references.bib
% =====================================================================

\section{Counterfactual Strategy for AI-READI SSMCGM-Stream}
\label{sec:aireadi_counterfactual_strategy}

\subsection{Motivation and scope}
\label{subsec:cf_motivation_scope}

The goal of the counterfactual component is to move SSMCGM-Stream from passive forecasting toward controlled, behavior-aware simulation. However, the AI-READI setting is fundamentally different from T1DEXI. T1DEXI contains explicit action streams such as timed insulin boluses, carbohydrate intake, and insulin-on-board, whereas AI-READI provides CGM, wearable signals, static clinical metadata, and a predicted meal flag. Therefore, AI-READI does not support insulin-dose or carbohydrate-dose counterfactuals. The variable \texttt{med\_insulin} is treated as a static participant descriptor, not as an editable action. All AI-READI counterfactuals are therefore framed as \emph{proxy scenario simulations}: meal-like, activity-like, stress-like, and sleep/rest physiological states.

The current SSMCGM-Stream run already demonstrates a strong participant-held-out forecaster, but its proxy scenario branch is not yet scientifically interpretable. Forecast-only and factual-future metrics are nearly identical, and the measured mean absolute scenario deltas are close to zero. This motivates a conservative strategy: before interpreting any scenario as a plausible counterfactual effect, we must first prove that the model uses the future scenario channel, that the edited scenarios lie on the empirical AI-READI manifold, and that model-implied effects agree with observed natural experiments.

The proposed counterfactual strategy has five linked components:
\begin{enumerate}
    \item a scenario-aware streaming forecasting model with explicit forecast-only and factual-future pathways;
    \item behavior-level, multivariate scenario edits rather than single-variable perturbations;
    \item trajectory-level support checks to enforce positivity and avoid out-of-distribution scenario paths;
    \item natural-experiment validation against matched observed windows;
    \item sensitivity analysis for unmeasured confounding, with additive bias bounds for continuous glucose effects and E-values for binary clinical risks.
\end{enumerate}
This section describes the mathematical and machine-learning framework behind these components.

\subsection{Streaming state-space forecasting backbone}
\label{subsec:cf_backbone}

Let participant $i$ have static covariates
\begin{equation}
    s_i = \{\mathrm{age}, \mathrm{HbA1c}, \mathrm{BMI}, \mathrm{site}, \mathrm{medications}, \ldots\}.
\end{equation}
These are encoded by a static encoder $g_\phi$ into an embedding
\begin{equation}
    e_i = g_\phi(s_i).
\end{equation}
At each 5-minute time step $t$, the dynamic observation vector $x_{i,t}$ contains observed CGM and wearable variables up to the current time, for example CGM, heart rate, steps, stress, and sleep stage. The streaming SSM update is
\begin{equation}
    h_{i,t} = f_\theta(h_{i,t-1}, x_{i,t}, e_i),
    \label{eq:stream_update}
\end{equation}
where $h_{i,t}$ is the participant-specific hidden state after observing the history up to time $t$. In the current SSMCGM-Stream implementation, $f_\theta$ is based on selective state-space/Mamba blocks. The practical advantage of this formulation is that the historical context is compressed into a fixed-dimensional state, so multiple future scenarios can be decoded from the same $h_{i,t}$ without re-encoding the full historical window.

For a forecast horizon of $H$ steps, the decoder receives three groups of future inputs:
\begin{itemize}
    \item $K_{t+1:t+H}$: known future inputs, such as time-of-day and calendar encodings;
    \item $A_{t+1:t+H}$: future scenario/proxy variables, such as meal proxy, activity footprint, stress, or sleep/rest signals;
    \item $M_{t+1:t+H}$: masks indicating whether scenario variables are known, unknown, factual, or edited.
\end{itemize}
The model predicts glucose quantiles through a residual-current target,
\begin{equation}
    \hat{y}^{(q)}_{i,t+h} = y_{i,t} + \hat{r}^{(q)}_{i,t,h}, \qquad h=1,\ldots,H,
    \label{eq:residual_current}
\end{equation}
where $y_{i,t}$ is the current observed glucose and $\hat{r}^{(q)}_{i,t,h}$ is the predicted residual quantile at horizon $h$. This makes persistence a strong baseline and prevents cold-start forecasts from predicting absolute glucose from scratch.

\subsection{Forecast-only and factual-future pathways}
\label{subsec:fo_ff_pathways}

The first requirement for counterfactual modeling is to check whether the scenario channel is alive. We therefore evaluate two decoder pathways for the same hidden state $h_{i,t}$.

The forecast-only pathway hides future scenario information:
\begin{equation}
    \hat{Y}^{fo}_{i,t+1:t+H}
    = D_\theta(h_{i,t}, K_{t+1:t+H}, A^0_{t+1:t+H}, M^0_{t+1:t+H}),
    \label{eq:fo_decoder}
\end{equation}
where $A^0$ is set to zero or neutral values and $M^0$ marks the future scenario path as unknown.

The factual-future pathway reveals observed future wearable/proxy variables retrospectively:
\begin{equation}
    \hat{Y}^{ff}_{i,t+1:t+H}
    = D_\theta(h_{i,t}, K_{t+1:t+H}, A^{obs}_{t+1:t+H}, M^{obs}_{t+1:t+H}).
    \label{eq:ff_decoder}
\end{equation}
This pathway is not deployable, because it uses information from the future. Its role is diagnostic: if the factual-future pathway does not improve prediction on event-rich windows, then the model has not learned to use scenario information.

Using the pinball loss $\rho_q(u)=\max(q u,(q-1)u)$, the multi-horizon quantile losses are
\begin{align}
    \mathcal{L}_{fo}
    &= \sum_{q \in \mathcal{Q}} \sum_{h=1}^{H}
    \rho_q\left(y_{i,t+h} - \hat{y}^{fo,(q)}_{i,t+h}\right), \\
    \mathcal{L}_{ff}
    &= \sum_{q \in \mathcal{Q}} \sum_{h=1}^{H}
    \rho_q\left(y_{i,t+h} - \hat{y}^{ff,(q)}_{i,t+h}\right).
\end{align}
The factual-future gain is
\begin{equation}
    G_{ff} = \mathcal{L}_{fo} - \mathcal{L}_{ff}.
    \label{eq:factual_gain}
\end{equation}
A useful scenario channel should satisfy $G_{ff}>0$ on event-rich anchors. If $G_{ff}\approx 0$, any inference-time scenario edit will mostly be arbitrary because the decoder was trained to ignore the future scenario path.

\subsection{Scenario-aware training objective}
\label{subsec:cf_training_objective}

To make the scenario path useful, the training objective should include both forecast-only and factual-future losses:
\begin{equation}
    \mathcal{L}_{train}
    = \mathcal{L}_{fo}
    + \lambda_{ff}\mathcal{L}_{ff}
    + \lambda_{gain}\mathcal{L}_{gain}
    + \lambda_{bal}\mathcal{L}_{bal}.
    \label{eq:scenario_train_loss}
\end{equation}
The gain term is applied only to event-rich anchors, where future wearable/proxy information is expected to matter:
\begin{equation}
    \mathcal{L}_{gain}
    = \mathbbm{1}_{event}
    \left[m + \mathcal{L}_{ff} - \mathcal{L}_{fo}\right]_+,
    \label{eq:gain_hinge}
\end{equation}
where $m \geq 0$ is a small margin. This loss penalizes the model when the factual-future path is not at least $m$ better than the forecast-only path. Importantly, we do \emph{not} include any term of the form $-\mathbb{E}|\Delta|$. Such a term would reward large scenario effects and could manufacture fake counterfactuals. The model is encouraged to use real future information when it is predictive, but it is never directly rewarded for making scenario effects large.

The optional balancing term $\mathcal{L}_{bal}$ is inspired by counterfactual recurrent networks. Let $B_t$ denote an observed behavior-proxy label, such as high activity, sleep/rest, stress, or meal-proxy. A behavior classifier $c_\psi$ attempts to predict $B_t$ from the hidden history representation $h_{i,t}$:
\begin{equation}
    \hat{B}_t = c_\psi(h_{i,t}).
\end{equation}
A domain-adversarial objective can be written as
\begin{equation}
    \min_\theta \max_\psi
    \left[\mathcal{L}_{train} - \lambda_{bal}\, CE(B_t,c_\psi(h_{i,t}))\right].
    \label{eq:adversarial_balancing}
\end{equation}
The purpose is to reduce the association between patient history and behavior assignment, which is a source of time-dependent confounding. This term should be introduced only after simpler two-pathway training and event oversampling have been tested, because excessive balancing can degrade forecasting performance.

\subsection{Behavior-level scenarios rather than symptom edits}
\label{subsec:behavior_level_scenarios}

A central design choice is to define scenarios at the level of behaviors, not individual wearable variables. Directly setting heart rate to a high value is not a valid intervention on exercise: heart rate is a physiological consequence of many possible causes, including activity, stress, illness, caffeine, and glucose dynamics. Therefore, AI-READI scenarios are represented by behavior labels
\begin{equation}
    B \in \{\mathrm{meal\_proxy},\mathrm{mild\_activity},\mathrm{high\_activity},
    \mathrm{sleep\_rest},\mathrm{sedentary\_stress}\}.
\end{equation}
Each behavior is realized as a multivariate wearable footprint:
\begin{equation}
    B \longrightarrow
    (\mathrm{steps}, \mathrm{heart\ rate}, \mathrm{stress}, \mathrm{sleep\ stage}, \mathrm{awake}, \ldots).
\end{equation}
This avoids impossible single-variable counterfactuals and is closer to a target-trial emulation view: the intervention is a behavioral strategy, and the wearable variables are its observed proxy signature.

\subsection{On-manifold empirical scenario path generation}
\label{subsec:on_manifold_generation}

A first attempt could set scenario variables using independent percentiles, for example steps at the participant's 75th percentile and heart rate at the 65th percentile. This is useful for prototyping, but it can create impossible combinations because marginally plausible values are not necessarily jointly plausible. The preferred approach is empirical conditional path sampling.

For each anchor $(i,t)$, define the matching context
\begin{equation}
    C_{i,t} = \left[
    y_{i,t-L:t},
    \Delta y_{i,t-L:t},
    HR_{i,t-L:t},
    steps_{i,t-L:t},
    stress_{i,t-L:t},
    sleep_{i,t-L:t},
    hour(t),
    s_i
    \right].
    \label{eq:matching_context}
\end{equation}
For a behavior $B$, the future scenario path is sampled from observed training windows with similar context:
\begin{equation}
    A^{B}_{i,t+1:t+H}
    \sim
    p_{train}\left(A_{u+1:u+H}\mid B_u=B, C_u \approx C_{i,t}\right).
    \label{eq:empirical_scenario_sampling}
\end{equation}
In practice, this can be implemented by nearest-neighbor retrieval, cluster prototypes, or weighted averaging among observed analogue windows. This makes the generated scenario path lie closer to the empirical AI-READI manifold than independently constructed percentiles.

The actual edit is localized in time and scope. For a behavior $B$ starting at horizon step $\tau$ and lasting $d$ steps, the scoped edit operator is
\begin{equation}
    (A^B,M^B,S^B)
    = \mathcal{E}_{B,\tau,d}(A,M,S;C_{i,t}),
    \label{eq:edit_operator}
\end{equation}
where $S$ is a source/type variable distinguishing unknown future values, factual future values, planned proxy values, and edited proxy scenarios. For affected variables $j\in\mathcal{V}_B$,
\begin{equation}
    A^B_{t+h,j}
    =
    \begin{cases}
        \tilde{A}^{B}_{h,j}, & h\in[\tau,\tau+d],\ j\in\mathcal{V}_B, \\
        A_{t+h,j}, & \mathrm{otherwise},
    \end{cases}
\end{equation}
where $\tilde{A}^{B}$ is the empirical analogue path. The mask is set to one for edited variables during the intervention interval, and the source/type label marks them as edited proxy values.

\subsection{Scenario predictions and model-implied effects}
\label{subsec:scenario_effect}

Given the same hidden state $h_{i,t}$, the forecast-only prediction is
\begin{equation}
    \hat{Y}^{fo}_{i,t+1:t+H}
    = D_\theta(h_{i,t},K,A^0,M^0,S^0),
\end{equation}
whereas the behavior-scenario prediction is
\begin{equation}
    \hat{Y}^{B}_{i,t+1:t+H}
    = D_\theta(h_{i,t},K,A^B,M^B,S^B).
\end{equation}
The model-implied proxy scenario effect at horizon $h$ is
\begin{equation}
    \Delta_\theta(B,i,t,h)
    = \hat{y}^{B}_{i,t+h} - \hat{y}^{fo}_{i,t+h}.
    \label{eq:model_effect}
\end{equation}
This is not automatically a causal effect. It is the change in the model's prediction when the future proxy path is changed while holding the observed history and current state fixed. It becomes scientifically interpretable only if the scenario channel is used, the scenario has empirical support, and the model-implied effect agrees with observed natural experiments.

A useful decomposition is
\begin{equation}
    \hat{y}^{B}_{i,t+h}
    = \hat{y}^{base}_{i,t+h} + \Delta_\theta(B,i,t,h),
    \label{eq:base_delta_decomposition}
\end{equation}
where $\hat{y}^{base}$ captures the forecast from history and known future time variables, and $\Delta_\theta$ captures the additional scenario-dependent component. This decomposition should be used diagnostically: if $\mathbb{E}|\Delta_\theta|\approx 0$, the scenario branch is inactive or weakly used.

\subsection{Trajectory-level support and positivity}
\label{subsec:support_positivity}

Longitudinal counterfactual identification requires positivity: each scenario should have nonzero probability among participants with comparable histories. In this application, this means the edited future wearable path must resemble observed AI-READI paths. Support should therefore be assessed over the whole trajectory rather than at a single time point.

Let $\Psi(A_{t+1:t+H},C_t)$ summarize a future path and its context, including total steps, maximum heart rate, mean stress, sleep fraction, time of onset, glucose level, glucose slope, time of day, site, HbA1c quartile, and study group. Define the path support distance
\begin{equation}
    d_{path}(B,i,t)
    =
    \min_{u\in\mathcal{D}^{B}_{train}}
    \left\|\Psi(A^{B}_{i,t+1:t+H},C_{i,t})
    - \Psi(A^{obs}_{u+1:u+H},C_u)\right\|_{\Sigma^{-1}}.
    \label{eq:path_support}
\end{equation}
The support class is then
\begin{equation}
    support(B,i,t)=
    \begin{cases}
        \mathrm{supported}, & d_{path}<c_1 \text{ and } n_{neighbors}>n_{min},\\
        \mathrm{weak}, & d_{path}<c_2,\\
        \mathrm{unsupported}, & \mathrm{otherwise}.
    \end{cases}
    \label{eq:support_class}
\end{equation}
Only supported scenarios are used in the main counterfactual analysis. Weakly supported scenarios are reported in the appendix. Unsupported scenarios are excluded from effect interpretation.

\subsection{Natural-experiment validation}
\label{subsec:natural_experiment_validation}

Because AI-READI is observational, model-implied effects must be compared against real observed windows that resemble the proposed scenario. For each behavior $B$, define a treatment indicator
\begin{equation}
    T^B_{i,t} = 1
\end{equation}
when the observed future window contains the behavior footprint, and $T^B_{i,t}=0$ for matched control windows. For example, an activity treated window may require high steps, elevated heart rate, and awake state; a control window should have low steps or normal heart rate while matching on pre-onset glucose and clinical context.

Matching should use rich pre-onset histories, not only current glucose and one slope. The matching vector is
\begin{equation}
    H^{match}_{i,t}=
    \left[y_{i,t-L:t},\Delta y_{i,t-L:t},HR_{i,t-L:t},steps_{i,t-L:t},stress_{i,t-L:t},sleep_{i,t-L:t},hour(t),s_i\right].
\end{equation}
The observed natural-experiment effect is
\begin{align}
    \Delta_{obs}(B,h)
    =
    &\ \mathbb{E}\left[y_{i,t+h}-y_{i,t}\mid T^B_{i,t}=1,\ matched\right]
    \\
    &-\mathbb{E}\left[y_{i,t+h}-y_{i,t}\mid T^B_{i,t}=0,\ matched\right].
    \label{eq:observed_effect}
\end{align}
The model-implied effect is averaged over the same supported anchors:
\begin{equation}
    \Delta_\theta(B,h)
    = \mathbb{E}_{(i,t)\in\mathcal{S}_B}
    \left[\hat{y}^{B}_{i,t+h}-\hat{y}^{fo}_{i,t+h}\right],
    \label{eq:average_model_effect}
\end{equation}
where $\mathcal{S}_B$ is the set of supported anchors for behavior $B$.

Validation metrics include effect MAE,
\begin{equation}
    EffectMAE(B)
    = \frac{1}{H}\sum_{h=1}^{H}
    \left|\Delta_\theta(B,h)-\Delta_{obs}(B,h)\right|,
\end{equation}
and sign agreement,
\begin{equation}
    SignAgreement(B)
    = \frac{1}{H}\sum_{h=1}^{H}
    \mathbbm{1}\left[\mathrm{sign}(\Delta_\theta(B,h))=\mathrm{sign}(\Delta_{obs}(B,h))\right].
\end{equation}
The primary validation report should include treated-window counts, control-window counts, support coverage, sign agreement, and effect MAE. Calibration slope is useful as a secondary metric, but it should not replace interpretable diagnostics.

\subsection{Sensitivity analysis for unmeasured confounding}
\label{subsec:sensitivity_analysis}

Sequential ignorability cannot be verified in AI-READI. Behavior may be confounded by unobserved intent, symptoms, meals, insulin, illness, or other latent factors. Sensitivity analysis is therefore a core component of the counterfactual framework rather than an optional add-on.

For continuous glucose effects in mg/dL, we use an additive bias bound. Let $\delta$ denote a hypothetical unmeasured confounding bias. The adjusted observed effect is
\begin{equation}
    \Delta_{obs}^{adj}(B,h;\delta)
    = \Delta_{obs}(B,h) - \delta.
\end{equation}
We evaluate a grid such as
\begin{equation}
    \delta \in \{-20,-15,-10,-5,0,5,10,15,20\}\ \mathrm{mg/dL}.
\end{equation}
The sign-flip threshold is
\begin{equation}
    \delta_{flip}(B,h)=
    \min_{\delta}\left\{\mathrm{sign}(\Delta_{obs}^{adj}(B,h;\delta))\neq
    \mathrm{sign}(\Delta_{obs}(B,h))\right\}.
\end{equation}
This gives a direct interpretation: how large an unmeasured glucose-scale bias would need to be to reverse the estimated effect.

For binary clinical outcomes, such as hypoglycemia within $H$ minutes,
\begin{equation}
    Y^{hypo}_{i,t,H}=\mathbbm{1}\left[\min_{1\leq h\leq H} y_{i,t+h}<70\right],
\end{equation}
or hyperglycemia,
\begin{equation}
    Y^{hyper}_{i,t,H}=\mathbbm{1}\left[\max_{1\leq h\leq H} y_{i,t+h}>180\right],
\end{equation}
we estimate risk ratios comparing treated and matched controls and report E-values. E-values are appropriate for risk-ratio-scale sensitivity analysis; they should not be used as the main sensitivity tool for continuous mg/dL effects unless those effects are converted into binary clinical risks.

\subsection{Attribution and channel-use diagnostics}
\label{subsec:attribution_diagnostics}

Attribution is useful only after the behavioral tests show that the scenario channel matters. Raw gradients such as
\begin{equation}
    \left\|\frac{\partial \hat{y}_{t+h}}{\partial A_{t+1:t+H}}\right\|
\end{equation}
can be unstable in gated SSMs. Therefore, the primary channel-use diagnostic remains the factual-future gain $G_{ff}$ on event windows. Attribution should be used as a secondary explanation tool.

When applied, attribution should answer three questions:
\begin{enumerate}
    \item Does the decoder place non-negligible attribution on the scenario path rather than only on current glucose?
    \item Does attribution increase after the planned scenario onset $\tau$?
    \item Does attribution vary by scenario and horizon in a physiologically plausible way?
\end{enumerate}
The hidden-attention interpretation of Mamba models can support this analysis by converting the selective SSM recurrence into an attention-like view of which past or scenario tokens influence the forecast.

\subsection{Counterfactual horizon design}
\label{subsec:horizon_design}

The current forecasting headline uses a 60-minute horizon, but some behavior effects may unfold over longer periods. Activity, stress, and sleep/rest effects can be delayed or sustained, and meal-like effects may peak beyond the first hour depending on timing. Therefore, the counterfactual validation should be run at several horizons:
\begin{equation}
    H \in \{12,24,48,72\}
\end{equation}
corresponding to 1, 2, 4, and 6 hours at 5-minute resolution. The 1-hour horizon remains the deployment-oriented forecast benchmark, while 2--6 hour horizons test whether behavioral scenario effects are hidden by the short forecasting window.

\subsection{Experimental roadmap}
\label{subsec:cf_experimental_roadmap}

The proposed implementation is divided into four phases.

\paragraph{Phase 0: audit the current model without retraining.}
Compute factual-future gain $G_{ff}$ and scenario deltas by all anchors and by event-rich anchors: meal-proxy windows, high-activity windows, sleep transitions, stress windows, and large glucose-change windows. If $G_{ff}\approx 0$, the current scenario decoder should not be interpreted.

\paragraph{Phase 1: train a scenario-aware stream model.}
Train variants with the two-pathway objective in Equation~\eqref{eq:scenario_train_loss}, event-enriched sampling, and optional source/type embeddings. The primary metric is not large scenario delta; it is positive factual-future gain on event windows.

\paragraph{Phase 2: replace hand-set scenarios with empirical on-manifold scenario paths.}
Compare independent percentile scenarios, nearest-neighbor empirical paths, and cluster-prototype paths. Report support coverage, path distance, effect stability, and scenario effect by horizon.

\paragraph{Phase 3: validate against natural experiments.}
For each behavior, compare model-implied effects with observed matched effects using sign agreement and effect MAE. Add continuous sensitivity bounds for mg/dL effects and E-values for binary hypo/hyper risk outcomes.

\paragraph{Phase 4: optional CRN-style balancing.}
If natural-experiment validation suggests substantial time-dependent confounding, add adversarial behavior balancing. This should be treated as an ablation because too much balancing may reduce forecast accuracy.

\subsection{Expected outputs}
\label{subsec:cf_expected_outputs}

The counterfactual pipeline should produce the following report-ready outputs:
\begin{itemize}
    \item \texttt{scenario\_channel\_audit.csv}: factual-future gain, event-window gain, and scenario delta by mode;
    \item \texttt{scenario\_support.csv}: support class, path distance, and neighbor counts by scenario;
    \item \texttt{scenario\_effects\_by\_horizon.csv}: $\Delta_\theta(B,h)$ by scenario and horizon;
    \item \texttt{natural\_experiment\_effects.csv}: treated/control counts, $\Delta_{obs}$, $\Delta_\theta$, sign agreement, and effect MAE;
    \item \texttt{sensitivity\_bounds.csv}: additive bias bounds for continuous effects and E-values for binary clinical risks;
    \item figures for factual-future gain, support coverage, model-vs-observed effects, and sensitivity thresholds.
\end{itemize}

\subsection{Interpretation rules}
\label{subsec:cf_interpretation_rules}

The final interpretation should follow strict rules:
\begin{enumerate}
    \item If $G_{ff}\leq 0$ on event windows, the scenario channel is considered inactive; do not interpret scenario deltas.
    \item If a scenario is unsupported by the trajectory-level support check, exclude it from the main counterfactual results.
    \item If model-implied effects disagree with observed natural-experiment effects, report the scenario as a failed or unvalidated proxy.
    \item If continuous effects are sensitive to small additive bias, state that the result is fragile to unmeasured confounding.
    \item E-values are reported only for binary clinical risk endpoints, not as the main sensitivity metric for continuous glucose effects.
    \item AI-READI counterfactuals should never be described as insulin-dose or carbohydrate-dose counterfactuals.
\end{enumerate}

Under these rules, the contribution is not a claim that AI-READI already supports definitive causal recommendations. The contribution is a validated framework for deciding which proxy behavioral scenarios are meaningful, which are unsupported, and which require additional action labels or model supervision.

